\documentclass[11pt,a4paper]{article}

% -------------------------------------------------------------------------
% Packages
% -------------------------------------------------------------------------
\usepackage[utf8]{inputenc}
\usepackage[T1]{fontenc}
\usepackage{amsmath,amssymb,amsthm}
\usepackage{booktabs}
\usepackage{graphicx}
\usepackage{hyperref}
\usepackage{xcolor}
\usepackage{microtype}
\usepackage{enumitem}
\usepackage{geometry}
\usepackage{natbib}
\usepackage{algorithm}
\usepackage{algpseudocode}
\usepackage{array}
\usepackage{multirow}

\geometry{margin=2.5cm}

% -------------------------------------------------------------------------
% Custom commands
% -------------------------------------------------------------------------
\newcommand{\themisdb}{\textsc{ThemisDB}}
\newcommand{\hypothesis}[1]{\textbf{Hypothesis:} \textit{#1}}
\newcommand{\caveat}[1]{\textcolor{gray}{\small\textit{Caveat: #1}}}
\newcommand{\TODO}[1]{\textcolor{red}{[TODO: #1]}}

\theoremstyle{definition}
\newtheorem{definition}{Definition}[section]
\newtheorem{proposition}{Proposition}[section]
\newtheorem{remark}{Remark}[section]

% -------------------------------------------------------------------------
% Title block
% -------------------------------------------------------------------------
\title{%
  \textbf{Boltzmann-Inspired Observability for AI Inference and RAG Pipelines:}\\
  \textbf{Energy-Entropy Monitoring, FLARE Parallels, and Correlation Analysis}\\[0.5em]
  \large\textit{ThemisDB Research Series --- ACTIVE DRAFT}
}

\author{%
  ThemisDB Research Group\\
  \texttt{research@themisdb.dev}\\[0.3em]
  \small Repository: \texttt{makr-code/ThemisDB}
}

\date{June 2026 --- Draft v0.1}

% -------------------------------------------------------------------------
\begin{document}
% -------------------------------------------------------------------------

\maketitle

\begin{abstract}
Monitoring the reliability of AI inference and Retrieval-Augmented Generation
(RAG) pipelines in production requires principled signal design beyond simple
latency and error-rate counters.
We propose a \emph{Boltzmann-inspired observability layer} that reframes
retrieval score distributions as a statistical-mechanical system, deriving
entropy, effective candidate count, and free-energy-like scalars as
first-class monitoring signals.
These signals provide interpretable, dimensionless proxies for retrieval
uncertainty, query-distribution drift, and potential hallucination risk.

We draw explicit parallels to the FLARE framework
\citep{jiang2023flare}, which triggers re-retrieval based on
token-level probability thresholds.
Our approach is complementary: whereas FLARE acts \emph{within} the
generation loop, the proposed layer acts \emph{across} the query log,
providing longitudinal drift and anomaly signals at the system level.

We describe how the approach integrates with \themisdb{}'s hybrid
graph-vector-relational architecture through concrete AQL query examples,
define an event-data schema for observability signals, and state
measurable performance expectations and SLO hypotheses for the monitoring
pipeline.
We clearly mark all empirical claims as hypotheses pending validation,
and include an explicit section on correlation-versus-causation caveats.
\end{abstract}

\paragraph{Keywords.}
Boltzmann distribution; entropy monitoring; RAG observability; FLARE;
distribution shift; ThemisDB; hybrid database; AQL; energy-based models;
uncertainty quantification; hallucination proxy.

\tableofcontents

\bigskip\hrule\bigskip

% =========================================================================
\section{Introduction}
% =========================================================================

Production AI systems built on RAG pipelines are subject to a class of
failure modes that are not well captured by traditional infrastructure
monitoring:
\begin{itemize}[noitemsep]
  \item \emph{Semantic drift} --- the query distribution shifts away from
    the domain covered by the indexed corpus.
  \item \emph{Retrieval uncertainty} --- the top-$k$ retrieved candidates
    concentrate on a single document (\emph{over-confident}) or are
    uniformly spread (\emph{under-determined}).
  \item \emph{Hallucination risk} --- the generator produces fluent text
    with little grounding support, often correlated with poor retrieval.
  \item \emph{Index staleness} --- the corpus diverges from the world
    modeled by incoming queries.
\end{itemize}

Standard observability stacks (Prometheus, OpenTelemetry) expose latency
histograms and error rates, but offer no semantic insight into \emph{why}
a RAG answer is unreliable.
Token-level approaches such as FLARE \citep{jiang2023flare} detect
low-confidence spans \emph{during} generation, but operate
single-query in real time and cannot accumulate longitudinal drift
signals across thousands of requests.

This paper proposes a complementary monitoring layer grounded in
statistical-mechanical reasoning:
by treating the retrieval score distribution as a Boltzmann/Gibbs
distribution, we derive dimensionless scalars --- entropy $H$, effective
candidate count $N_\text{eff}$, and a free-energy-like monitoring scalar
$\Phi$ --- that are interpretable, cheap to compute, and
amenable to rolling-window and cross-request analysis.

The contribution is deliberately \emph{architectural and methodological},
not empirical: we propose a design, state measurable hypotheses, and
define an evaluation protocol.
Empirical validation is deferred to future work.

\paragraph{Paper organization.}
Section~\ref{sec:background} reviews the statistical-mechanical and
information-theoretic foundations.
Section~\ref{sec:related} surveys related work.
Section~\ref{sec:method} defines the method in detail.
Section~\ref{sec:integration} describes the \themisdb{} integration.
Section~\ref{sec:performance} states performance expectations.
Section~\ref{sec:evaluation} defines the evaluation protocol.
Section~\ref{sec:threats} discusses threats to validity.
Section~\ref{sec:conclusion} concludes.

% =========================================================================
\section{Background}
\label{sec:background}
% =========================================================================

\subsection{The Boltzmann/Gibbs View of Probability}

In statistical mechanics, the \emph{Boltzmann distribution} assigns a
probability to a microstate $i$ with energy $E_i$ as:
\begin{equation}
  p_i = \frac{e^{-E_i / (k_B T)}}{Z},
  \qquad Z = \sum_j e^{-E_j / (k_B T)},
  \label{eq:boltzmann}
\end{equation}
where $T$ is temperature and $k_B$ is Boltzmann's constant.
By the \emph{maximum entropy principle} of Jaynes
\citep{jaynes1957}, this is the unique distribution that maximises
entropy subject to a constraint on the expected energy.

In machine learning, the \emph{softmax} function is mathematically
identical to Eq.~\eqref{eq:boltzmann} with $k_B T \equiv 1$ and
energies defined as negative logits.
This connection motivates treating neural network output distributions
as Gibbs distributions at a given effective temperature, an interpretation
exploited in energy-based models \citep{lecun2006,liu2020ebm} and
uncertainty-calibration research \citep{guo2017calibration}.

\subsection{Entropy, Uncertainty, and the Effective Candidate Count}

Given a probability distribution $\mathbf{p} = (p_1,\ldots,p_k)$,
Shannon entropy \citep{shannon1948} is:
\begin{equation}
  H(\mathbf{p}) = -\sum_{i=1}^{k} p_i \ln p_i.
  \label{eq:entropy}
\end{equation}
$H$ is zero for a point mass and $\ln k$ for the uniform distribution.

A useful derived quantity is the \emph{effective candidate count}:
\begin{equation}
  N_\text{eff}(\mathbf{p}) = e^{H(\mathbf{p})},
  \label{eq:neff}
\end{equation}
which equals 1 for a degenerate (fully concentrated) distribution and
$k$ for a uniform one.
$N_\text{eff}$ has an intuitive interpretation: how many candidates
does the system ``effectively consider''?
A value near 1 indicates over-confidence; a value near $k$ indicates
retrieval uncertainty.

\subsection{RAG Observability: Motivating the Problem}

Standard RAG pipelines (see \citet{lewis2020rag,gao2023ragsurvey})
combine a retriever (dense or sparse) with a generator LLM.
The retriever returns top-$k$ candidates scored by semantic similarity
(dot-product or cosine in embedding space).
These scores are typically consumed by the generator as context without
further monitoring.

Failure modes in production include:
\begin{itemize}[noitemsep]
  \item Single-document overreliance (low $N_\text{eff}$, high confidence
    in potentially wrong source).
  \item Retrieval failure (all scores near zero, high $H$ near $\ln k$,
    generator hallucinating).
  \item Distribution shift (entropy baseline drifts over time due to
    evolving query semantics or stale index).
\end{itemize}
None of these are visible to standard infrastructure monitors.

% =========================================================================
\section{Related Work}
\label{sec:related}
% =========================================================================

\subsection{Energy-Based Models and OOD Detection}

Energy-based models (EBMs) \citep{lecun2006} define a scalar energy
function $E_\theta(x)$ and derive probabilities via Eq.~\eqref{eq:boltzmann}.
\citet{liu2020ebm} show that the energy score is a strong baseline for
out-of-distribution (OOD) detection, outperforming softmax confidence.
Their free-energy formulation:
\begin{equation}
  F(x) = -\log \sum_{c} e^{f_c(x)},
  \label{eq:free_energy_ood}
\end{equation}
inspired our free-energy-like monitoring scalar (Section~\ref{sec:fscalar}).

\subsection{Uncertainty Quantification and Calibration}

\citet{guo2017calibration} demonstrate that modern neural networks are
miscalibrated: confidence scores do not accurately reflect empirical accuracy.
Temperature scaling (adjusting the effective $T$ in
Eq.~\eqref{eq:boltzmann}) is the simplest calibration method.
\citet{lakshminarayanan2017} and \citet{malinin2018predictive} provide
deeper uncertainty estimates via ensembles and prior networks, respectively.

\subsection{Distribution Shift and Concept Drift}

The concept-drift literature \citep{gama2014survey} provides
statistical tests (ADWIN, DDM, Page-Hinkley) for detecting
distributional changes in streaming data.
These methods are complementary to our energy-entropy approach:
they test for drift in scalar-valued time series, while we propose
what signals to put in those time series.

\subsection{RAG Reliability and Hallucination}

\citet{ji2023survey} survey hallucination in NLG.
\citet{min2023factscore} and \citet{es2023ragas} provide
factual-precision and end-to-end RAG evaluation frameworks (FActScore,
RAGAS).
These rely on ground-truth labels, which are expensive.
Our approach provides cheaper, label-free monitoring proxies.

\subsection{FLARE: Forward-Looking Active Retrieval}
\label{sec:flare_review}

FLARE \citep{jiang2023flare} is the closest related work.
During generation, FLARE monitors token probabilities.
When the predicted probability of the next token falls below a threshold
$\tau$, FLARE triggers a re-retrieval step using a reformulated query
derived from the current generation prefix.
The key design choices are:
\begin{itemize}[noitemsep]
  \item \textbf{Trigger signal:} minimum token probability $\min_t p(t)$
    or perplexity of the upcoming span.
  \item \textbf{Scope:} single query, real-time, within the generation
    loop.
  \item \textbf{Action:} re-retrieve and regenerate the low-confidence
    span.
\end{itemize}

Our approach shares the intuition that low confidence warrants
attention, but differs in scope (system-level, cross-query) and
granularity (retrieval scores rather than generation token probabilities).
Table~\ref{tab:flare_comparison} in Section~\ref{sec:flare_parallel}
provides a structured comparison.

% =========================================================================
\section{Method}
\label{sec:method}
% =========================================================================

\subsection{Boltzmannized Retrieval Energy Model}

Let $q$ be a query, and let $\mathcal{D}_k = \{d_1,\ldots,d_k\}$ be
the top-$k$ retrieved documents with raw similarity scores
$s_i = \text{sim}(q, d_i) \in \mathbb{R}$.
We define the \emph{retrieval energy} of candidate $i$ as:
\begin{equation}
  E_i(q) = -\alpha \cdot s_i - \beta \cdot r_i + \gamma \cdot \rho_i,
  \label{eq:energy}
\end{equation}
where:
\begin{itemize}[noitemsep]
  \item $\alpha > 0$: weight for semantic similarity (lower energy = more
    relevant).
  \item $r_i \in [0,1]$: recency score (normalized document age); higher
    recency reduces energy.
  \item $\rho_i \geq 0$: risk flag (e.g., low-quality source, negative
    feedback label); increases energy.
  \item $\beta, \gamma \geq 0$: tunable weights.
\end{itemize}

The simplest configuration with $\alpha=1$, $\beta=\gamma=0$ recovers
$E_i = -s_i$, making Boltzmann weights a direct re-normalization of
similarity scores.
Complex configurations allow incorporation of provenance quality and
temporal relevance.

We apply the Boltzmann normalization with inverse temperature $\beta_T = 1/T$:
\begin{equation}
  p_i = \frac{e^{-\beta_T E_i}}{\sum_{j=1}^{k} e^{-\beta_T E_j}},
  \qquad Z = \sum_{j=1}^{k} e^{-\beta_T E_j}.
  \label{eq:gibbs}
\end{equation}
For $T=1$ this is a standard softmax over negated energies.
The temperature $T$ acts as a smoothing parameter:
$T \to 0$ concentrates mass on the minimum-energy candidate;
$T \to \infty$ yields uniform $p_i = 1/k$.

\begin{remark}
  In production we do \emph{not} claim to solve a physical Boltzmann
  equation.
  The term ``Boltzmann-inspired'' denotes the use of this mathematical
  form for its well-understood statistical properties (maximum-entropy,
  calibrated probabilities), not a physical simulation.
\end{remark}

\subsection{Entropy, Effective Candidate Count, and Per-Query Statistics}

From $\mathbf{p}$ we derive the per-query monitoring signals:

\paragraph{Retrieval entropy.}
$H_q = H(\mathbf{p})$ as in Eq.~\eqref{eq:entropy}.
High $H_q$ near $\ln k$ indicates diffuse retrieval (many nearly
equally relevant or equally irrelevant candidates).
Low $H_q$ near $0$ indicates concentrated retrieval.

\paragraph{Effective candidate count.}
$N_{\text{eff},q} = e^{H_q}$ as in Eq.~\eqref{eq:neff}.
Interpretable: ``how many documents effectively determine the answer?''

\paragraph{Energy gap.}
$\Delta E_q = E_{(2)} - E_{(1)}$, the difference between the
second-lowest and lowest energy.
A small gap means the top-1 and top-2 candidates are nearly
indistinguishable; a large gap means the system has a clear winner.

\paragraph{Partition function proxy.}
$\ln Z_q = \ln \sum_j e^{-\beta_T E_j}$.
Unusual values (very low $Z$) indicate that all candidates have high
energy, suggesting out-of-distribution queries.

\subsection{Free-Energy-Like Monitoring Scalar}
\label{sec:fscalar}

Inspired by the Helmholtz free energy $F = U - TS$ and the EBM
formulation of \citet{liu2020ebm}, we define a scalar
\emph{monitoring proxy}:
\begin{equation}
  \Phi_q = \bar{E}_q - T \cdot H_q
         = \sum_i p_i E_i + T \sum_i p_i \ln p_i,
  \label{eq:fscalar}
\end{equation}
where $\bar{E}_q = \sum_i p_i E_i$ is the expected energy under $\mathbf{p}$.

$\Phi_q$ is low when retrieval is confident (low energy) and high-entropy
(many relevant candidates), and high when retrieval is uncertain or
concentrated on high-energy (low-quality) candidates.
It serves as a single-number summary of retrieval health per query.

\subsection{Correlation Analysis Design}
\label{sec:correlation}

We define three primary correlation hypotheses to be tested empirically:

\begin{enumerate}
  \item \textbf{Quality vs. Entropy.}
    \hypothesis{$H_q$ is negatively correlated with user-labeled
    relevance scores or RAGAS context-precision.}
    Intuition: concentrated retrieval (low $H$) should correlate
    with higher precision when the top candidate is correct.
    \caveat{The relationship may be non-monotone: very low $H$ may
    indicate over-confident but wrong retrieval.}

  \item \textbf{Hallucination proxy vs. Energy gap.}
    \hypothesis{Small $\Delta E_q$ is positively associated with
    hallucination-flagged responses (LLM self-consistency or
    FActScore $< \tau$).}
    Intuition: when the top-2 candidates are nearly equivalent in
    energy, the generator lacks a clear grounding anchor.
    \caveat{Ground-truth hallucination labels are expensive; surrogate
    labels (self-consistency, model-based NLI) introduce noise.}

  \item \textbf{Latency vs. Entropy.}
    \hypothesis{High $H_q$ is associated with increased end-to-end
    latency due to longer context (more candidates passed to generator)
    or re-retrieval triggered downstream.}
    \caveat{Confounders include query complexity, generator load, and
    index size.}
\end{enumerate}

\paragraph{Correlation vs. causation caveat (mandatory).}
All of the above are observational hypotheses.
Even if positive correlations are found, they do not establish causal
mechanisms.
For example, a positive correlation between $H_q$ and hallucination
rate may reflect a common cause (queries about rare topics) rather
than entropy causing hallucination.
Causal analysis would require controlled interventions (e.g., ablating
retrieval diversity) which are outside the scope of this proposal.

\subsection{FLARE Parallels and Differences}
\label{sec:flare_parallel}

Table~\ref{tab:flare_comparison} summarises the structural comparison
between FLARE and the proposed Boltzmann observability layer.

\begin{table}[ht]
  \centering
  \caption{FLARE vs. Boltzmann Observability Layer: structural comparison.}
  \label{tab:flare_comparison}
  \begin{tabular}{@{}lll@{}}
    \toprule
    \textbf{Dimension} & \textbf{FLARE} & \textbf{Boltzmann Layer (proposed)} \\
    \midrule
    Scope & Single query, real-time & Cross-query, longitudinal \\
    Signal source & Generator token probabilities & Retriever score distributions \\
    Primary metric & $\min_t p(t_i)$ / perplexity & $H_q$, $N_{\text{eff},q}$, $\Phi_q$ \\
    Action & Re-retrieve \& regenerate & Alert / drift report / AQL query \\
    Loop integration & Inside generation loop & Outside, post-hoc or async \\
    Latency impact & Direct (adds retrieval calls) & Negligible (async logging) \\
    Temporal analysis & No (single request) & Yes (rolling windows, baselines) \\
    Hallucination response & Proactive (prevents) & Reactive (detects pattern) \\
    Ground-truth needed & No & No (proxy signals) \\
    Corpus coverage & Implicitly via re-retrieval & Explicitly via drift metrics \\
    Implementation locus & LLM inference service & Database observability layer \\
    ThemisDB integration & Via RAG module hooks & Native AQL + relational store \\
    \bottomrule
  \end{tabular}
\end{table}

The two approaches are complementary: FLARE corrects individual
low-confidence generations; the Boltzmann layer detects systemic
patterns (drift, systematic overconfidence, index coverage gaps) that
a per-request approach cannot expose.

% =========================================================================
\section{ThemisDB Architecture Integration}
\label{sec:integration}
% =========================================================================

\themisdb{} is a hybrid graph-vector-relational database with an
AQL query interface and MVCC transaction semantics
\citep{themisdb2026}.
The hybrid model makes it unusually well-suited for observability
workloads that mix structured event records (relational), embedding
comparisons (vector), and provenance tracking (graph).

\subsection{Data Model (Event Schema)}

Each RAG pipeline invocation emits an observability event stored in
\themisdb{}:

\begin{verbatim}
ObservabilityEvent {
  event_id       : UUID          -- primary key
  timestamp      : TIMESTAMP     -- wall-clock, UTC
  query_hash     : BYTES(32)     -- SHA-256 of query embedding
  session_id     : UUID          -- optional; links to user session
  topk_scores    : FLOAT64[]     -- raw similarity scores s_1..s_k
  topk_doc_ids   : UUID[]        -- retrieved document identifiers
  energy_params  : JSON          -- alpha, beta, gamma, T used
  -- Derived signals (computed on ingest):
  entropy_H      : FLOAT64       -- H(p), see Eq. (3)
  n_eff          : FLOAT64       -- exp(H), see Eq. (4)
  energy_gap     : FLOAT64       -- E_(2) - E_(1)
  ln_Z           : FLOAT64       -- log partition function
  phi_scalar     : FLOAT64       -- free-energy proxy, see Eq. (6)
  -- Optional downstream signals:
  answer_confidence : FLOAT64    -- surrogate confidence (if available)
  latency_ms     : INT64         -- end-to-end request latency
  feedback_label : SMALLINT      -- -1/0/1 user feedback (nullable)
}
\end{verbatim}

\subsection{Mapping to ThemisDB Storage Model}

\begin{itemize}[noitemsep]
  \item \textbf{Relational store:} \texttt{ObservabilityEvent} table,
    rolling baselines per time bucket, SLO alert records.
  \item \textbf{Vector store:} query embeddings indexed for similarity
    search, enabling ``find queries most similar to this anomalous one''.
  \item \textbf{Graph store:} provenance graph linking event nodes to
    document nodes, document nodes to source nodes.
    Enables attribution analysis: ``which source documents are most
    often associated with high-$\Phi$ (risky) retrievals?''
  \item \textbf{File / blob store:} long-term archival of raw events;
    sampling frames for offline evaluation.
\end{itemize}

\subsection{Pipeline Stages}

\begin{enumerate}[label=\textbf{Stage \arabic*.}, leftmargin=*]
  \item \textbf{Ingest.} RAG service posts \texttt{ObservabilityEvent}
    to \themisdb{} via gRPC / HTTP write API.
    Derived signals ($H$, $N_\text{eff}$, $\Phi$) computed server-side
    on ingest to avoid client-side implementation inconsistency.
  \item \textbf{Rolling Baseline.} Materialized views maintain
    rolling $\mu$, $\sigma$, median, MAD of $H_q$, $\Phi_q$
    per 1h / 24h / 7d windows.
    Baseline is recomputed incrementally using sliding-window aggregation.
  \item \textbf{Alerting.} Rule engine evaluates:
    \begin{itemize}[noitemsep]
      \item $|H_q - \mu_H| > z \cdot \text{MAD}_H$ (entropy spike),
      \item $\Phi_q > \Phi_{95\%}$ of baseline (free-energy anomaly),
      \item $\frac{1}{W}\sum_{w} \mathbf{1}[\Phi_q > \theta]
            > \alpha_\text{alert}$ (sustained elevated risk).
    \end{itemize}
    Alerts stored in \texttt{AlertEvent} table with reasoning payload.
  \item \textbf{Analysis.} AQL queries (Section~\ref{sec:aql_examples})
    enable operator investigation of flagged events, cluster analysis,
    and correlation queries.
\end{enumerate}

\subsection{AQL Query Examples}
\label{sec:aql_examples}

\paragraph{Entropy drift (time-series view).}
\begin{verbatim}
SELECT
  DATE_TRUNC(timestamp, 'hour') AS bucket,
  AVG(entropy_H) AS mean_entropy,
  PERCENTILE_CONT(0.95) WITHIN GROUP (ORDER BY entropy_H) AS p95_entropy,
  COUNT(*) AS n
FROM ObservabilityEvent
WHERE timestamp >= NOW() - INTERVAL '7 days'
GROUP BY bucket
ORDER BY bucket;
\end{verbatim}

\paragraph{High-risk low-entropy events (over-confident failures).}
\begin{verbatim}
SELECT event_id, timestamp, entropy_H, phi_scalar, feedback_label
FROM ObservabilityEvent
WHERE entropy_H < 0.5          -- concentrated retrieval
  AND phi_scalar > :phi_p95    -- high free-energy (risky)
  AND feedback_label = -1      -- negative user feedback
ORDER BY phi_scalar DESC
LIMIT 100;
\end{verbatim}

\paragraph{Domain-shift cluster candidates.}
\begin{verbatim}
SELECT
  o.event_id,
  o.entropy_H,
  o.phi_scalar,
  v.cluster_id          -- from vector k-means materialized view
FROM ObservabilityEvent o
JOIN QueryCluster v ON v.query_hash = o.query_hash
WHERE o.timestamp >= NOW() - INTERVAL '1 day'
  AND v.cluster_id NOT IN (
    SELECT DISTINCT cluster_id FROM QueryCluster
    WHERE snapshot_date < NOW() - INTERVAL '30 days'
  )
ORDER BY o.phi_scalar DESC;
\end{verbatim}

% =========================================================================
\section{Performance Expectations and SLO Hypotheses}
\label{sec:performance}
% =========================================================================

This section states \emph{hypotheses} about acceptable performance
targets for the observability layer.
These are engineering targets, not yet empirically validated results.

\subsection{Ingestion Throughput}

\hypothesis{The observability ingest pipeline sustains
$\geq 5{,}000$ events/second on \themisdb{}'s standard write path
without degrading primary query throughput by more than 2\%.}

Basis: derived signal computation ($H$, $N_\text{eff}$, $\Phi$) requires
$O(k)$ floating-point operations per event with $k \leq 100$.
At $k=20$ and 64-bit floats, this is approximately $\sim 1$~\textmu s
CPU per event, negligible versus I/O.

\subsection{Query Latency Targets}

\hypothesis{The following p50/p95/p99 latency targets hold for
key monitoring queries under a 30-day rolling event store
($\sim 10^8$ events):}

\begin{table}[ht]
  \centering
  \caption{Query latency SLO hypotheses.}
  \label{tab:latency}
  \begin{tabular}{@{}llrrr@{}}
    \toprule
    Query type & Index strategy & p50 & p95 & p99 \\
    \midrule
    Hourly entropy drift (7d) & Time-partition & $<50$~ms & $<200$~ms & $<500$~ms \\
    High-risk events lookup & $(\Phi, H)$ composite index & $<20$~ms & $<100$~ms & $<300$~ms \\
    Domain-shift cluster join & Vector + relational & $<100$~ms & $<400$~ms & $<1$~s \\
    Per-session trace (1 session) & session\_id index & $<10$~ms & $<50$~ms & $<100$~ms \\
    \bottomrule
  \end{tabular}
\end{table}

\caveat{These targets assume the event table has columnar compression
on numeric signals and appropriate partition pruning.
Without index tuning, p99 may exceed targets by an order of magnitude.}

\subsection{Storage Overhead}

\hypothesis{The observability event store adds $\leq 10\%$ of total
\themisdb{} storage for a typical RAG deployment processing
$10^6$ queries/day with $k=20$ candidates.}

Estimation:
$\sim 400$~bytes per event (UUID + timestamps + $k$ floats + derived
signals + metadata) $\times$ $10^6$/day $= 400$~MB/day uncompressed.
With columnar compression ($\sim$4$\times$), $\approx 100$~MB/day.

\subsection{Alert Precision/Recall Hypotheses}

\hypothesis{A simple threshold alert on $\Phi_q > \Phi_{95\%}$
achieves $\geq 0.6$ precision and $\geq 0.5$ recall for
hallucination-flagged requests, measured against FActScore $< 0.5$
as ground truth.}

\caveat{Precision/recall are sensitive to the base rate of
hallucinations, which varies widely by domain and model.
This hypothesis should be tested on domain-specific evaluation sets.
The 0.6/0.5 values are \emph{minimum acceptable} targets, not
guaranteed results.}

\subsection{Computational Complexity}

\begin{itemize}[noitemsep]
  \item Per-event signal computation: $O(k)$ time, $O(k)$ space.
  \item Rolling window baseline update (sliding sum): $O(1)$ amortized
    with ring-buffer implementation.
  \item Top-$k$ anomaly detection: $O(n \log k)$ with priority queue,
    where $n$ is events in the window.
  \item Domain-shift cluster candidates query: $O(n \cdot d)$ for
    vector similarity where $d$ is embedding dimension;
    HNSW index reduces to $O(\log n)$ expected for ANN retrieval.
\end{itemize}

% =========================================================================
\section{Evaluation Protocol}
\label{sec:evaluation}
% =========================================================================

\subsection{Benchmark Plan}

\paragraph{Synthetic workload.}
Generate $10^5$ synthetic queries with controlled entropy levels
($H \in \{0.1, 0.5, 1.0, \ln 20\}$) by constructing score vectors with
known distributions.
Measure: ingest throughput, signal computation accuracy (vs. reference
Python implementation), alert trigger latency.

\paragraph{Replay workload.}
Replay a 30-day event log from a staging RAG deployment.
Inject synthetic anomalies (entropy spikes, domain-shift batches) at
known timestamps.
Measure: detection latency, false-positive rate over anomaly-free periods.

\paragraph{End-to-end RAG evaluation.}
Use a labeled dataset (e.g., Natural Questions, TriviaQA, or internal
\themisdb{} test sets) with known hallucination labels.
Measure: Pearson/Spearman correlation of $H_q$, $\Phi_q$, $\Delta E_q$
with FActScore and RAGAS context-precision.

\paragraph{Correlation analysis.}
For each hypothesis in Section~\ref{sec:correlation}:
\begin{enumerate}[noitemsep]
  \item Report Pearson $r$ and Spearman $\rho$ with 95\% bootstrap
    confidence intervals.
  \item Report partial correlations controlling for query length and
    $k$.
  \item Report effect size (Cohen's $d$ for binary outcomes).
\end{enumerate}

% =========================================================================
\section{Threats to Validity}
\label{sec:threats}
% =========================================================================

\paragraph{Internal validity.}
Signal computation assumes that raw similarity scores are comparable
across queries.
In practice, scores depend on the embedding model and normalization.
Cross-model or cross-index comparisons require re-calibration of
energy parameters $(\alpha, \beta, \gamma, T)$.

\paragraph{External validity.}
Results from one RAG deployment may not generalise to others with
different domain, document quality, or query types.
The evaluation protocol uses multiple datasets to partially mitigate this.

\paragraph{Construct validity.}
The $\Phi_q$ scalar is a heuristic proxy for ``retrieval health'',
not a rigorously derived measure.
Its correlation with ground-truth quality metrics must be established
per deployment before use in automated alerting.

\paragraph{Correlation vs. causation.}
As noted in Section~\ref{sec:correlation}, all observed correlations
between entropy/energy signals and quality outcomes are observational.
No causal claims are made in this paper.
Confounders (query difficulty, topic rarity, model bias) are not
controlled for without randomised experimental design.

\paragraph{Temperature sensitivity.}
The choice of $T$ in Eq.~\eqref{eq:gibbs} affects all derived signals.
If $T$ is poorly calibrated, signals will be compressed (high $T$)
or degenerate (low $T$).
A calibration procedure (e.g., Platt scaling on a validation set)
is required before deployment.

% =========================================================================
\section{Conclusion}
\label{sec:conclusion}
% =========================================================================

We have proposed a Boltzmann-inspired observability layer for RAG pipelines
and AI inference, grounded in the statistical-mechanical reinterpretation
of retrieval score distributions.
The core signals --- entropy $H$, effective candidate count $N_\text{eff}$,
energy gap $\Delta E$, and free-energy proxy $\Phi$ --- are cheap to
compute, interpretable, and amenable to rolling-window drift analysis.

We drew explicit parallels to FLARE, showing that the two approaches
are complementary: FLARE acts within the generation loop on individual
requests; the Boltzmann layer acts across the query log at the system
level.
The proposed integration with \themisdb{}'s hybrid storage model
enables AQL-driven correlation queries, provenance graph analysis,
and multi-tier retention policies in a single system.

Performance expectations and evaluation protocols are stated as
testable hypotheses.
The next step is empirical validation on representative RAG deployments,
measurement of signal-to-quality correlations, and iterative calibration
of the energy model parameters.

\paragraph{Future work.}
(1)~Causal analysis of entropy-quality relationships via controlled
re-retrieval ablations.
(2)~Online learning of energy parameters from user feedback.
(3)~Extension to multi-modal retrieval (text + image + structured data).
(4)~Integration with \themisdb{}'s MVCC to provide temporally
consistent snapshots of observability baselines.

\bigskip

\paragraph{Acknowledgements.}
This draft builds on the existing \themisdb{} RAG module, Shannon
information theory survey
(\texttt{research/PAPER\_SURVEY\_SHANNON\_LLM\_RAG\_FLARE.md}),
and the community of contributors to the open ThemisDB repository.

% =========================================================================
% References
% =========================================================================
\bibliographystyle{plainnat}
\bibliography{bib/boltzmann_flare_rag_monitoring}

\end{document}
